\documentclass[a4paper,12pt]{article}
\usepackage[utf8]{inputenc}
\usepackage[T1]{fontenc}
\usepackage[french]{babel}
\usepackage{amsmath, amssymb, amsfonts}
\usepackage{array, longtable, booktabs}
\usepackage[table]{xcolor} % option table pour \rowcolors
\usepackage{geometry}
\geometry{left=1.8cm, right=1.8cm, top=2cm, bottom=2cm}
\usepackage{enumitem}
\usepackage{hyperref}
\hypersetup{colorlinks=true, linkcolor=blue, urlcolor=blue}

% Commandes pour les coches
\newcommand{\fait}{\ensuremath{\checkmark}}
\newcommand{\revoit}{\ensuremath{\circlearrowright}}
\newcommand{\casevide}{\ensuremath{\square}}

% Types de colonnes pour un contrôle précis des largeurs
\newcolumntype{C}[1]{>{\centering\arraybackslash}p{#1}}
\newcolumntype{L}[1]{>{\raggedright\arraybackslash}p{#1}}
\newcolumntype{R}[1]{>{\raggedleft\arraybackslash}p{#1}}

\title{Programmation annuelle de Mathématiques expertes -- Terminale générale \\
	Année 2026--2027 -- Zone C (3 h/semaine)}
\author{}
\date{Version du \today}

\begin{document}
	
	\maketitle
	
	\section*{Présentation générale}
	Ce document propose une organisation spiralaire de l’enseignement de l’option Mathématiques expertes, conforme au programme officiel. L’année comporte 31 semaines de cours effectifs (hors vacances et jours fériés). Chaque notion est introduite une première fois, puis réinvestie et approfondie dans un second temps, à l’occasion de thèmes d’étude différents. Cette approche favorise une acquisition progressive et durable des connaissances.
	
	\medskip
	\textbf{Place de Python :} La programmation est intégrée de façon transverse et évaluée dans \emph{tous} les sujets (formatives, contrôles, DM, PBL). Les compétences algorithmiques (boucles, tests, fonctions, simulation, calcul matriciel, calcul approché) sont développées en parallèle des mathématiques et réactivées régulièrement selon le même principe spiralé.
	
	\medskip
	\textbf{Utilisation de l’IA dans les DM :} Dans chaque devoir maison, l’élève est invité à utiliser un outil d’intelligence artificielle générative (ChatGPT, Mistral, Copilot, etc.) pour :
	\begin{itemize}
		\item obtenir des explications ou des exemples supplémentaires sur une notion,
		\item suggérer des pistes de résolution,
		\item débuguer ou améliorer un code Python,
		\item confronter différentes méthodes de résolution.
	\end{itemize}
	\emph{L’IA ne doit cependant pas rédiger la solution finale à la place de l’élève.} Le raisonnement personnel et la rédaction restent exigibles. Un \textbf{journal de bord} (questions posées, réponses obtenues, critiques et adaptations apportées) est à annexer à chaque DM. Cette pratique développe l’esprit critique et la capacité à évaluer la pertinence des réponses générées.
	
	\newpage
	
	\section{Tableau récapitulatif des contenus à enseigner}
	Le tableau ci-dessous reprend l’intégralité des attendus du programme. La colonne \textbf{« Traité »} indique si le contenu a été abordé en classe (coche \fait), et la colonne \textbf{« Reprise »} signale les notions qui feront l’objet d’un second passage (\revoit) pour consolider la compréhension.
	
	\rowcolors{2}{gray!8}{white}
	\begin{longtable}{|L{14cm}|C{1.5cm}|C{1.5cm}|}
		\hline
		\textbf{Thème / Contenu} & \textbf{Traité} & \textbf{Reprise} \\
		\hline
		\endhead
		\hline
		\multicolumn{3}{r}{Suite de la page suivante} \\
		\endfoot
		\hline
		\endlastfoot
		
		% --- Nombres complexes – Algèbre ---
		\multicolumn{3}{|L{14cm}|}{\textbf{Nombres complexes – Algèbre}} \\
		Ensemble $\mathbb{C}$, partie réelle, imaginaire, opérations & \casevide & \casevide \\
		Conjugaison, propriétés algébriques & \casevide & \casevide \\
		Inverse d’un complexe non nul & \casevide & \casevide \\
		Formule du binôme dans $\mathbb{C}$ & \casevide & \casevide \\
		Résolution d’équation linéaire $az=b$ & \casevide & \casevide \\
		Résolution d’équation simple avec $z$ et $\bar{z}$ & \casevide & \casevide \\
		\hline
		
		% --- Nombres complexes – Géométrie ---
		\multicolumn{3}{|L{14cm}|}{\textbf{Nombres complexes – Géométrie}} \\
		Image d’un complexe, affixe d’un point/vecteur & \casevide & \casevide \\
		Module, interprétation géométrique, $|z|^2=z\bar{z}$ & \casevide & \casevide \\
		Module d’un produit, d’un inverse & \casevide & \casevide \\
		Ensemble $\mathbb{U}$ des complexes de module 1 & \casevide & \casevide \\
		Arguments, forme trigonométrique & \casevide & \casevide \\
		\hline
		
		% --- Nombres complexes et trigonométrie ---
		\multicolumn{3}{|L{14cm}|}{\textbf{Complexes et trigonométrie}} \\
		Formules d’addition/duplication (produit scalaire) & \casevide & \casevide \\
		Exponentielle imaginaire, notation $e^{i\theta}$, forme exponentielle & \casevide & \casevide \\
		Formules d’Euler, formule de Moivre & \casevide & \casevide \\
		\hline
		
		% --- Équations polynomiales ---
		\multicolumn{3}{|L{14cm}|}{\textbf{Équations polynomiales}} \\
		Solutions complexes d’une équation du second degré à coeff. réels & \casevide & \casevide \\
		Factorisation $z^n-a^n$ par $z-a$ & \casevide & \casevide \\
		Si $P(a)=0$, factorisation par $z-a$ & \casevide & \casevide \\
		Un polynôme de degré $n$ a au plus $n$ racines & \casevide & \casevide \\
		\hline
		
		% --- Utilisation des complexes en géométrie ---
		\multicolumn{3}{|L{14cm}|}{\textbf{Complexes en géométrie}} \\
		Interprétation géométrique du module et de l’argument & \casevide & \casevide \\
		Racines $n$-ièmes de l’unité, ensemble $\mathbb{U}_n$ & \casevide & \casevide \\
		\hline
		
		% --- Arithmétique ---
		\multicolumn{3}{|L{14cm}|}{\textbf{Arithmétique}} \\
		Divisibilité dans $\mathbb{Z}$ & \casevide & \casevide \\
		Division euclidienne, congruences, compatibilité & \casevide & \casevide \\
		PGCD, algorithme d’Euclide & \casevide & \casevide \\
		Couples d’entiers premiers entre eux & \casevide & \casevide \\
		Théorème de Bézout & \casevide & \casevide \\
		Théorème de Gauss & \casevide & \casevide \\
		Nombres premiers, infinité, décomposition en facteurs premiers & \casevide & \casevide \\
		Petit théorème de Fermat & \casevide & \casevide \\
		\hline
		
		% --- Graphes et matrices ---
		\multicolumn{3}{|L{14cm}|}{\textbf{Graphes et matrices}} \\
		Graphe, sommets, arêtes, degré, ordre, chaîne, connexité & \casevide & \casevide \\
		Matrice (tableau), carrée, colonne, ligne, opérations & \casevide & \casevide \\
		Inverse, puissances d’une matrice carrée & \casevide & \casevide \\
		Représentations : matrice d’adjacence, transformations géométriques, systèmes, suites & \casevide & \casevide \\
		Chaînes de Markov à 2 ou 3 états, matrice de transition, distribution initiale & \casevide & \casevide \\
		Interprétation des coefficients de $P^n$, distribution après $n$ transitions & \casevide & \casevide \\
		Distributions invariantes & \casevide & \casevide \\
		\hline
	\end{longtable}
	
	\section{Calendrier et programmation harmonieuse}
	
	\subsection{Calendrier retenu (zone C, jours ouvrés)}
	Les vacances et jours fériés ont été pris en compte. Les semaines de cours sont les suivantes :
	
	\begin{center}
		\begin{tabular}{|C{3.5cm}|C{5cm}|C{2.5cm}|}
			\hline
			\textbf{Période} & \textbf{Dates} & \textbf{Nb de semaines} \\
			\hline
			Rentrée – Toussaint & 1er sept. – 16 oct. 2026 & 7 \\
			Toussaint – Noël & 2 nov. – 18 déc. 2026 & 7 \\
			Noël – Hiver & 4 janv. – 5 fév. 2027 & 5 \\
			Hiver – Printemps & 22 fév. – 2 avr. 2027 & 6 \\
			Printemps – Fin mai & 19 avr. – 31 mai 2027 & 6 \\
			\hline
			\textbf{Total} & & \textbf{31 semaines} \\
			\hline
		\end{tabular}
	\end{center}
	
	\textit{Jours fériés (tombant un jour de cours) exclus :} 1er janvier 2027 (vendredi), 29 mars 2027 (lundi de Pâques), 6 mai 2027 (Ascension, jeudi), 17 mai 2027 (Pentecôte, lundi). Les 1er et 8 mai 2027 tombent un samedi.
	
	\subsection{Programmation par trimestre (approche spiralée)}
	Les tableaux suivants précisent pour chaque semaine :
	\begin{itemize}
		\item les contenus nouveaux,
		\item les \textbf{points ciblés dans l’évaluation formative} (5 à 10 min en début de séance) incluant systématiquement une question Python,
		\item des remarques utiles.
	\end{itemize}
	Les évaluations formatives mêlent mathématiques et programmation, avec réactivation des notions antérieures.
	
	% --- 1er trimestre ---
	\newpage
	\subsubsection{1\textsuperscript{er} trimestre (septembre – novembre 2026) – 14 semaines}
	
	\rowcolors{2}{gray!8}{white}
	\begin{longtable}{|C{0.8cm}|L{4.0cm}|L{5.8cm}|L{2.5cm}|}
		\hline
		\textbf{Sem.} & \textbf{Contenus nouveaux} & \textbf{Évaluation formative -- points évalués (dont Python)} & \textbf{Remarques} \\
		\hline
		\endhead
		\hline
		\multicolumn{4}{r}{Suite de la page suivante} \\
		\endfoot
		\hline
		\endlastfoot
		
		1-2 & Nombres complexes : définition, opérations, conjugué, inverse. & Calculs algébriques sur les complexes. \textbf{Python} : fonction qui calcule le conjugué et l’inverse. & Introduction algébrique. \\
		\hline
		3-4 & Module, argument, forme trigonométrique, $\mathbb{U}$. & \textbf{Spirale :} opérations sur les complexes (S1-2). \textbf{Python} : module et argument avec \texttt{cmath}. & Géométrie. \\
		\hline
		5-6 & Formules d’Euler, de Moivre, exponentielle complexe. & \textbf{Spirale :} module/argument (S3-4). \textbf{Python} : calcul de puissances et tracé de polygones. & Trigonométrie. \\
		\hline
		7-8 & Équations du second degré à coeff. réels. Factorisation $z^n-a^n$. & \textbf{Spirale :} forme exponentielle (S5-6). \textbf{Python} : résolution d’équations polynomiales. & Polynômes. \\
		\hline
		9-10 & Arithmétique : divisibilité, division euclidienne, congruences. & \textbf{Spirale :} factorisation (S7-8). \textbf{Python} : test de divisibilité, calcul de PGCD. & Congruences. \\
		\hline
		11-12 & PGCD, algorithme d’Euclide, Bézout, Gauss. & \textbf{Spirale :} congruences (S9-10). \textbf{Python} : algorithme d’Euclide étendu. & Théorèmes fondamentaux. \\
		\hline
		13-14 & Nombres premiers, infinité, décomposition, petit théorème de Fermat. & \textbf{Bilan T1 :} complexes, arithmétique. \textbf{Python} : crible d’Ératosthène. & Primalité. \\
		\hline
		\multicolumn{4}{l}{} \\
		\hline
		\multicolumn{4}{l}{\textbf{Contrôles continus :} semaines 5, 10, 14 (avec exercice Python). \textbf{Formatives :} chaque semaine.} \\
		\hline
	\end{longtable}
	
	% --- 2e trimestre ---
	\newpage
	\subsubsection{2\textsuperscript{e} trimestre (novembre 2026 – février 2027) – 12 semaines}
	
	\rowcolors{2}{gray!8}{white}
	\begin{longtable}{|C{0.8cm}|L{4.0cm}|L{5.8cm}|L{2.5cm}|}
		\hline
		\textbf{Sem.} & \textbf{Contenus nouveaux} & \textbf{Évaluation formative -- points évalués (dont Python)} & \textbf{Remarques} \\
		\hline
		\endhead
		\hline
		\multicolumn{4}{r}{Suite de la page suivante} \\
		\endfoot
		\hline
		\endlastfoot
		
		15-16 & Graphes : définitions, adjacence, degré, chaîne, connexité. & \textbf{Spirale :} complexes (S1-8). \textbf{Python} : représentation d’un graphe par liste d’adjacence. & Modélisation. \\
		\hline
		17-18 & Matrices : définitions, opérations, inverse, puissances. & \textbf{Spirale :} graphes (S15-16). \textbf{Python} : multiplication matricielle, calcul de puissances. & Calcul matriciel. \\
		\hline
		19-20 & Matrice d’adjacence, nombre de chemins. Transformations géométriques. & \textbf{Spirale :} matrices (S17-18). \textbf{Python} : calcul de $A^n$ et interprétation. & \textbf{PBL 1 (2h)} : Marche aléatoire sur un graphe. \\
		\hline
		21-22 & Systèmes linéaires, suites récurrentes avec matrices. & \textbf{Spirale :} arithmétique (S9-14). \textbf{Python} : résolution de systèmes avec \texttt{numpy}. & Applications. \\
		\hline
		23-24 & Chaînes de Markov : états, matrice de transition, distribution initiale. & \textbf{Spirale :} matrices (S17-18). \textbf{Python} : simulation d’une chaîne de Markov. & Probabilités. \\
		\hline
		25-26 & Distributions invariantes, interprétation de $P^n$. & \textbf{Spirale :} Markov (S23-24). \textbf{Python} : calcul de la distribution stationnaire. & \textbf{PBL 2 (2h)} : Modèle de diffusion d’Ehrenfest. \\
		\hline
		\multicolumn{4}{l}{} \\
		\hline
		\multicolumn{4}{l}{\textbf{Contrôles continus :} semaines 18, 22, 26 (avec exercice Python). \textbf{Formatives :} chaque semaine.} \\
		\hline
	\end{longtable}
	
	% --- 3e trimestre ---
	\newpage
	\subsubsection{3\textsuperscript{e} trimestre (février – mai 2027) – 12 semaines}
	
	\rowcolors{2}{gray!8}{white}
	\begin{longtable}{|C{0.8cm}|L{4.0cm}|L{5.8cm}|L{2.5cm}|}
		\hline
		\textbf{Sem.} & \textbf{Contenus nouveaux} & \textbf{Évaluation formative -- points évalués (dont Python)} & \textbf{Remarques} \\
		\hline
		\endhead
		\hline
		\multicolumn{4}{r}{Suite de la page suivante} \\
		\endfoot
		\hline
		\endlastfoot
		
		27-28 & Racines $n$-ièmes de l’unité, polygones réguliers. & \textbf{Spirale :} complexes (S1-8). \textbf{Python} : tracé des racines avec \texttt{matplotlib}. & Applications géométriques. \\
		\hline
		29-30 & Équations de degré 3 à coeff. réels (une racine connue). Formules de Viète. & \textbf{Spirale :} polynômes (S7-8). \textbf{Python} : factorisation automatique. & \textbf{PBL 3 (2h)} : Pentagone régulier – construction et calcul. \\
		\hline
		31-32 & Utilisation des complexes en géométrie (alignement, orthogonalité, etc.). & \textbf{Spirale :} racines (S27-28). \textbf{Python} : calculs de distances, angles. & Géométrie analytique. \\
		\hline
		33-34 & Arithmétique approfondie : équations diophantiennes, RSA. & \textbf{Spirale :} arithmétique (S9-14). \textbf{Python} : implémentation du chiffrement RSA. & Cryptographie. \\
		\hline
		35-36 & Chaînes de Markov : applications (PageRank, etc.). & \textbf{Bilan matrices/Markov}. \textbf{Python} : simulation de PageRank. & Synthèse. \\
		\hline
		37-38 & Révisions et évaluations finales. & \textbf{Bilan général :} tous les thèmes + Python. & Préparation épreuves. \\
		\hline
		\multicolumn{4}{l}{} \\
		\hline
		\multicolumn{4}{l}{\textbf{Contrôles continus :} semaines 30, 34, 38 (avec exercice Python). \textbf{Formatives :} chaque semaine.} \\
		\hline
	\end{longtable}
	
	\newpage
	
	\section{Évaluations}
	
	\subsection{Évaluations formatives (hebdomadaires)}
	Chaque semaine, une courte activité (5–10 min) mêle mathématiques et Python. Les points évalués sont détaillés dans les tableaux ci-dessus. Ces évaluations ne sont pas notées, mais elles permettent une auto-évaluation et un suivi des compétences algorithmiques.
	
	\subsection{Évaluations sommatives (contrôles continus)}
	Au moins trois contrôles par trimestre, d’une durée de 1 h chacun.  
	Chaque contrôle comporte :
	\begin{itemize}
		\item une partie sans calculatrice (calculs, raisonnement),
		\item une partie avec calculatrice / Python (exercice de programmation : écriture de fonction, complétion de code, simulation, résolution numérique, etc.).
	\end{itemize}
	
	\begin{center}
		\begin{tabular}{|C{3.5cm}|C{3.5cm}|C{3.5cm}|C{3.5cm}|}
			\hline
			\textbf{Trimestre} & \textbf{Contrôle 1} & \textbf{Contrôle 2} & \textbf{Contrôle 3} \\
			\hline
			1\textsuperscript{er} (sem. 5, 10, 14) & Complexes (algèbre, trigo) \newline + Python (calculs) & Équations polynomiales \newline + Python (résolution) & Arithmétique (PGCD, congruences) \newline + Python (Euclide) \\
			\hline
			2\textsuperscript{e} (sem. 18, 22, 26) & Graphes, matrices (opérations) \newline + Python (puissances) & Systèmes linéaires, suites \newline + Python (numpy) & Chaînes de Markov \newline + Python (simulation) \\
			\hline
			3\textsuperscript{e} (sem. 30, 34, 38) & Racines de l’unité, géométrie complexe \newline + Python (tracé) & RSA, équations diophantiennes \newline + Python (cryptographie) & Synthèse (tous thèmes) \newline + Python (projet complet) \\
			\hline
		\end{tabular}
	\end{center}
	
	\subsection{Problèmes basés sur l’apprentissage (PBL) – 2 h chacun}
	Trois PBL sont répartis sur l’année, en séances de 2 h (travail en groupes, compte rendu écrit et code Python rendu) :
	
	\begin{enumerate}
		\item \textbf{PBL 1 (semaine 20)} : \emph{Marche aléatoire sur un graphe}. Simuler une marche aléatoire, calculer les probabilités de transition à l’aide de la matrice d’adjacence, comparer avec les résultats théoriques.
		\item \textbf{PBL 2 (semaine 26)} : \emph{Modèle de diffusion d’Ehrenfest}. Modéliser un échange de molécules entre deux urnes par une chaîne de Markov, déterminer la distribution invariante, simuler avec Python.
		\item \textbf{PBL 3 (semaine 30)} : \emph{Construction du pentagone régulier}. Utiliser les racines $5$-ièmes de l’unité pour construire un pentagone, calculer les lignes trigonométriques, illustrer avec Python.
	\end{enumerate}
	
	\section{Devoirs maison (DM) -- progression, Python et utilisation de l’IA}
	
	Afin de favoriser une acquisition à long terme, un devoir maison est distribué \textbf{toutes les deux semaines} (18 DM sur l’année).  
	Chaque DM porte sur l’ensemble du programme traité depuis le 1er septembre, et comporte :
	\begin{itemize}
		\item des exercices de calcul et de raisonnement,
		\item un exercice de programmation Python,
		\item une mission spécifique d’utilisation de l’IA (consultation, débuggage, confrontation de méthodes ou critique de réponse).
	\end{itemize}
	
	L’élève doit joindre à sa copie un \textbf{journal de bord} retraçant ses échanges avec l’IA et les décisions prises à la suite de ces échanges.
	
	\begin{center}
		\small
		\rowcolors{2}{gray!8}{white}
		\begin{longtable}{|C{0.8cm}|C{1cm}|L{6.8cm}|L{6.8cm}|}
			\hline
			\textbf{DM} & \textbf{Sem.} & \textbf{Thèmes évalués (mathématiques + Python)} & \textbf{Mission IA spécifique} \\
			\hline
			\endhead
			\hline
			\multicolumn{4}{r}{Suite à la page suivante} \\
			\endfoot
			\hline
			\endlastfoot
			
			1 & 3 & Complexes : algèbre, conjugué, inverse. Python : calculs. & Demander à l’IA d’expliquer la notion de conjugué et de l’illustrer avec un exemple. \\
			\hline
			2 & 5 & Module, argument, forme trigonométrique. Python : calcul de module/argument. & Utiliser l’IA pour générer un nombre complexe aléatoire et en donner les formes. \\
			\hline
			3 & 7 & Formules d’Euler/de Moivre. Python : tracé de polygones. & Demander à l’IA de démontrer la formule de Moivre par récurrence, puis vérifier par Python. \\
			\hline
			4 & 9 & Équations du second degré à coeff. réels. Python : résolution automatique. & Soumettre à l’IA une équation et comparer les méthodes de résolution. \\
			\hline
			5 & 11 & Arithmétique : divisibilité, congruences. Python : test de divisibilité. & Demander à l’IA d’expliquer le test de divisibilité par 9 et de l’illustrer. \\
			\hline
			6 & 13 & PGCD, Bézout. Python : algorithme d’Euclide. & Utiliser l’IA pour trouver une combinaison linéaire de deux nombres. \\
			\hline
			7 & 16 & Graphes : définition, adjacence. Python : création d’un graphe. & Demander à l’IA de proposer un exemple de graphe modélisant un réseau social. \\
			\hline
			8 & 18 & Matrices : opérations, inverse. Python : multiplication. & Utiliser l’IA pour vérifier l’inverse d’une matrice et interpréter les erreurs. \\
			\hline
			9 & 20 & Puissances de matrices, nombre de chemins. Python : calcul de $A^n$. & Demander à l’IA d’interpréter les coefficients de $A^3$ dans un graphe. \\
			\hline
			10 & 22 & Systèmes linéaires. Python : résolution avec numpy. & Utiliser l’IA pour générer un système et comparer les méthodes de résolution. \\
			\hline
			11 & 24 & Chaînes de Markov : matrice de transition. Python : simulation. & Demander à l’IA de simuler une chaîne à 2 états et de vérifier la stationnarité. \\
			\hline
			12 & 26 & Distribution invariante. Python : calcul de la stationnaire. & Utiliser l’IA pour résoudre le système linéaire associé et discuter de l’unicité. \\
			\hline
			13 & 28 & Racines de l’unité. Python : tracé. & Demander à l’IA d’expliquer la construction du pentagone régulier. \\
			\hline
			14 & 30 & Équations de degré 3. Python : factorisation. & Utiliser l’IA pour factoriser un polynôme et vérifier la racine. \\
			\hline
			15 & 32 & Géométrie complexe (alignement, orthogonalité). Python : calculs. & Demander à l’IA de démontrer une propriété géométrique avec les complexes. \\
			\hline
			16 & 34 & RSA : chiffrement/déchiffrement. Python : implémentation. & Utiliser l’IA pour générer une paire de clés RSA et expliquer les étapes. \\
			\hline
			17 & 36 & PageRank. Python : simulation. & Demander à l’IA d’expliquer l’algorithme PageRank et de l’appliquer à un petit graphe. \\
			\hline
			18 & 38 & Bilan général. Python : projet complet. & Demander à l’IA un retour sur la structure du code et proposer des améliorations. \\
			\hline
		\end{longtable}
	\end{center}
	
	\textbf{Remarques :}
	\begin{itemize}
		\item Les DM sont à rendre une semaine après la distribution.
		\item La mission IA doit être clairement identifiée dans le journal de bord (copie d’écran ou retranscription des questions/réponses, suivie de l’analyse personnelle de l’élève).
		\item La correction est effectuée en classe, avec une attention particulière aux erreurs fréquentes, aux méthodes de résolution et aux bonnes pratiques d’utilisation de l’IA.
	\end{itemize}
	
	\section{Conseils pour la mise en œuvre}
	\begin{itemize}
		\item \textbf{Rituels} : consacrer 5 à 10 minutes par séance à des exercices de calcul (mental, littéral) et à de courts extraits de code Python (lecture, complétion, débogage).
		\item \textbf{Travail en îlots} : favoriser les échanges entre élèves lors des phases de recherche et des PBL.
		\item \textbf{Numérique} : utiliser Python régulièrement (au moins une séance sur deux). Insister sur la programmation modulaire (fonctions) et la documentation du code.
		\item \textbf{IA éthique} : former les élèves à formuler des requêtes précises et à évaluer de manière critique les réponses de l’IA (vérification des calculs, cohérence, limites). Insister sur le fait que l’IA est un \emph{outil d’aide}, pas un substitut au raisonnement personnel.
		\item \textbf{Différenciation} : proposer des exercices de difficulté progressive et des approfondissements pour les élèves plus à l’aise (notamment en programmation et en interaction avec l’IA).
		\item \textbf{Trace écrite} : veiller à ce que chaque élève dispose d’un cahier de cours structuré (définitions, propriétés, exemples) et d’un cahier de programmation regroupant les fonctions et scripts clés.
	\end{itemize}
	
	\section*{Conclusion}
	Cette programmation respecte les 31 semaines de cours disponibles et couvre l’intégralité du programme de mathématiques expertes, avec une intégration systématique de Python et de l’utilisation critique de l’IA dans les apprentissages et les évaluations. La progressivité spiralaire, appliquée aussi aux compétences algorithmiques et numériques, garantit que les notions mathématiques, la programmation et le jugement critique sont solidement maîtrisés en fin d’année. Les évaluations variées (formatives, sommatives, PBL, DM) permettent de suivre la progression de chaque élève et de développer ses compétences de modélisation, de raisonnement, de communication, de codage et de collaboration avec les outils d’intelligence artificielle.
	
\end{document}